%!TEX root = ./../../springer-mono-format/book.tex
%!TEX root = ./../../thesis-trier.tex
 
\IN{strongly regular codes}
Solid codes provide outstanding fault-tolerance when used for information transmission through a noisy channel w.r.t.~symbol substitutions and synchronization errors. In this context,
they were introduced by Levenshtein~\cite{Lev64} under the name \emph{strongly regular codes}, 
and further investigated by Romanov~\cite{Rom66}.
Shyr and Yu~\cite{ShyrYu90} introduced solid codes to investigate disjunctive domains, 
but did not recognize the connection to the earlier work by Levenshtein and Romanov. 
Later, Jürgensen and Yu~\cite{DBLP:journals/eik/JurgensenY90} observed that strongly 
regular codes describe precisely the solid codes.
For further information, we refer to the surveys~\cite{DBLP:journals/iandc/Jurgensen08,DBLP:reference/hfl/JurgensenK97},
or the introductory section of~\cite{DBLP:journals/tcs/JurgensenS21}. 
Recall that \cref{constrsync:ex:binarysolidcodeitem,constrsync:ex:constr-solid-code}
from \cref{constrsync:ex:binarysolidcode} gave an infinite binary solid code 
and a method to construct a solid codes, respectively.

In our investigation on the constrained synchronization problem and synchronization conditions
for automata obtained in reductions, we discovered (to the best of our knowledge)
new characterizations of solid codes, which we present in the following
(see the discussion section,
\cref{constr_sync:sec:gen_results:discussion}, page~\pageref{discussion:strongly-self-sync-codes}, 
for some background on how this characterization arose in the context
of the constrained synchronization problem).

Here, first, in \cref{thm:char-solid-codes}, we recap known characterizations. 
\Cref{thm:solidcode} collects further characterizations of solid codes
which are new and connect solid codes to various language equations.
Lastly, \cref{thm:solidcode-char2,cor:solidcode-char3,constrsync:prop:sigma-solid-code-char} 
give two other new characterizations.

\begin{theorem} \label{thm:char-solid-codes}
Let $L \subseteq \Sigma^+$ with $L \ne \emptyset$. Then the following are equivalent.
\begin{enumerate}
	
\item  \label{thm:char-solid-codes:1} $L$ is a solid code.
	
\item  \label{thm:char-solid-codes:2} (Jürgensen \& Yu~\cite{DBLP:journals/eik/JurgensenY90}) For every word $w \in \Sigma^*$ 
 we have $w = u_0 w_1 u_1 w_2 u_2 \cdots w_n u_n$
 for some $n \ge 0$ and unique words $u_0, u_1, \ldots, u_n \in \Sigma^* \setminus \Sigma^* L \Sigma^*$
 and $w_1, w_2, \ldots, w_n \in L$.
	
\item \label{thm:char-solid-codes:3} (Staiger~\cite{DBLP:journals/tcs/JurgensenS21}, Romanov~\cite{Rom66}, Levenshtein~\cite{Lev64})
 There exists a state-invariant decoder without look-ahead for $L$ (see \cite[Def.~9]{DBLP:journals/tcs/JurgensenS21}
 for a definition of these decoders for finite solid codes and extensions to infinite solid codes).

\item \label{thm:char-solid-codes:4} (Liu, Guo \& Shum~\cite[Thm.~2.4]{Liu2021}) $L$ is an infix code and $\{ u_1\cdots u_n \}$ is overlap-free for 
 all sequences of distinct words $u_1, \ldots, u_n \in L$ with $n \ge 1$.\footnote{In~\cite[Thm.~2.4]{Liu2021},
 	if~$L$ is finite, then overlap-freeness is only stated to hold for $\{u_1\cdots u_n\}$
 	if $2 \le n \le |L|$. However, this seems to be a typo,
 	as for example $L = \{ aba, bcb \}$
 	is an infix code
 	and $\{ ababcb\}$ and $\{ bcbaba \}$
 	are each overlap-free, while~$L$
 	is not overlap-free and so not a solid code.} 
 
\item \label{thm:char-solid-codes:5} (Liu, Guo \& Shum~\cite[Thm.~2.4]{Liu2021}) $L$ is an infix code such that $\{u\}$ is overlap-free for all $u \in L$ and $\{uv\}$ is overlap-free for $u,v \in L$ with $u \ne v$.
\end{enumerate}	
\end{theorem}

\begin{remark}
	Since for distinct $u_1, \ldots, u_n \in L$
	we have $u_i \cdots u_n \ne u_1 \cdots u_{n - i + 1}$
	for all $i \in \{1,\ldots,n\}$.
	Thus~$L$ is overlap-free
	if $\{ u_1 \cdots u_n \}$
	is overlap-free for distinct $u_1, \ldots, u_n \in L$.
	Next, if for distinct $u,v \in L$
	we have $u = xy$ and $v = yz$
	for non-empty $x,y,z \in \Sigma^+$,
	then $vu = yzxy$ and so $\{vu\}$
	is not overlap-free.
	Thus, if $\{u\}$ for $u \in L$ and $\{uv\}$
	for distinct $u,v \in L$ are overlap-free,
	then~$L$ is overlap-free.
	This shows that the additional conditions in Items~\ref{thm:char-solid-codes:4} and~\ref{thm:char-solid-codes:5} of \cref{thm:char-solid-codes}, beyond being an infix code, are equivalent to overlap-freeness (without any additional assumptions on~$L$).
\end{remark}

\begin{theorem} \label{thm:solidcode}
Let $L \subseteq \Sigma^+$. Then the following are equivalent:
\begin{enumerate}
\item  \label{thm:solidcode:1} $L$ is a solid code.

\item  \label{thm:solidcode:1a} $L$ is a bifix code and $(\pref{L} \setminus L) L \cap \Sigma^* L \Sigma^+ = \emptyset$.

\item  \label{thm:solidcode:1b}  $L$ is a suffix code and $\left( \bigcup_{u \in L} ( \pref{u}\setminus\{u\} ) \right)L \cap \Sigma^*L\Sigma^+ = \emptyset$.

\item \label{thm:solidcode:2} $L$ is an infix code and $(\pref{L} \setminus L) L \cap L \Sigma^+ = \emptyset$.

\item \label{thm:solidcode:3} $L$ is an infix code and $\pref{L} \cap \suff{L} = L \cup \{\varepsilon\}$.	
\end{enumerate}
\end{theorem}
\begin{proof}
$\ref{thm:solidcode:1}\Rightarrow \ref{thm:solidcode:1a}$: A solid code is an infix code, and so a bifix code, and overlap-free.
Suppose $vw = xuy$
with $x \in \Sigma^*$, $y \in \Sigma^+$, $u, w \in L$, and
$v \in \pref{L} \setminus L$.
By \cref{lem:overlap-free-prefix}
we have $|x| > 0$.
For $|x| \ge |v|$ (resp.~$|y| \ge |w|$), we have $u \in \factor{w}$ (resp.~$u \in \factor{v}$).
For $0 < |x| < |v|$ and $0 < |y| < |w|$,
we have $u = v'w'$ for a non-empty
and proper suffix $v'$ (resp.~prefix~$w'$)
of $v$ (resp.~$w$), so that $L$ is not overlap-free.
Hence, for $(\pref{L} \setminus L) L \cap \Sigma^+ L \Sigma^+  \ne \emptyset$,
we get that $L$ is not a solid code, which gives the claim.

$\ref{thm:solidcode:1a}\Rightarrow \ref{thm:solidcode:1b}$: Since~$L$ is a prefix code,
 its set of proper prefixes equals $\pref{L}\setminus L$, which gives the claim.

$\ref{thm:solidcode:1b}\Rightarrow \ref{thm:solidcode:2}$: Suppose~$L$ is a suffix code as stated.
Assume we have $uvw \in L$ with $u \in \Sigma^*$, $v \in L$,
and $w \in \Sigma^+$.
Let $x$ be a proper prefix of some word in~$L$.
Then $xuvw \in \Sigma^* L \Sigma^+$
and this is excluded by assumption.
Thus, $L$ is an infix code.
The assumption yields that $( \pref{L}\setminus L ) L \cap L \Sigma^+ = \emptyset$.

$\ref{thm:solidcode:2} \Rightarrow \ref{thm:solidcode:3}$: 	By \cref{lem:overlap-free-prefix}, $L$ is overlap-free.
Thus, \cref{constrsync:prop:overlap-free-inclusion-eqn}
gives $\pref{L} \cap \suff{L} \subseteq L \cup \{\varepsilon\}$.
The inclusion
$L \cup \{\varepsilon\} \subseteq \pref{L} \cap \suff{L}$
holds for any language.	

$\ref{thm:solidcode:3} \Rightarrow \ref{thm:solidcode:1}$: 
By \cref{constrsync:prop:pref-suff-overlap-free-inclusion-eqn}, $L$
is overlap-free and thus a solid code.	
\end{proof}

\begin{remark}
	As far as we are aware of,
	the condition $(\pref{L} \setminus L) L \cap \Sigma^* L \Sigma^+ = \emptyset$
	was never investigated in the literature before (apart
	from our own work~\cite{DBLP:conf/fct/Hoffmann21} 
	whose results are presented and extended
	in the present thesis).
	Note that $\pref{L}L \subseteq L\Sigma^+$ holds for every language $L \subseteq \Sigma^*$.
\end{remark}

We mention an implication for languages  $L \subseteq \Sigma^*$
such that $(\pref{L} \setminus L) L \cap \Sigma^* L \Sigma^+ = \emptyset$.

\begin{proposition}
	Let $L \subseteq \Sigma^*$
	such that $(\pref{L} \setminus L) L \cap \Sigma^* L \Sigma^+ = \emptyset$.
	Then $\pref{L}\subseteq L$
	or $uvw \in L$ with $u,v \in \Sigma^*$ and $v \in L$ implies $w = \varepsilon$.
\end{proposition}
\begin{proof}
	Suppose we have  $x \in \pref{L} \setminus L$
	for some $x \in \Sigma^*$.
	Let $v \in L$
	with $uvw \in L$ for some $u,w \in \Sigma^*$.
	Assume $|w| > 0$.
	Then we have $xuvw \in (\pref{L}\setminus L) L \cap \Sigma^* L \Sigma^+$.
	Thus, we must have $|w| = 0$.
\end{proof}

\begin{theorem} \label{thm:solidcode-char2}
	Let $L \subseteq \Sigma^+$.
	Then $L$ is a solid code
	if and only if for all $u \in \pref{L}$ and $v \in L$
	such that $uv = xwy$ for  $x \in \Sigma^*$ and $y \in \Sigma^+$ and $w \in L$
	we have $x = \varepsilon$ and $u = w$.			
\end{theorem}
\begin{proof}
	First, suppose $L$ is a solid code.
	Let $u \in \pref{L}$, $v, w \in L$,
	and $uv = xwy$ for $x \in \Sigma^*$ and $y \in \Sigma^+$.	
	By \crefitem{thm:solidcode}{thm:solidcode:1a}, we must have $u \in L$.
	By \crefitem{constr_sync:prop:selfsync_implies_avoiding}{constr_sync:prop:item:sc-in-cf}, 
	$L$ is comma-free.
	So $u, v \in L$	imply $x = \varepsilon$ and, as $L$ is a prefix code, $w = u$.	
	Conversely, suppose the stated condition holds.
	Then, clearly, $L$ is comma-free.
	Next, suppose $u \in \pref{L} \setminus L$ and $v,w \in L$
	with $uv = xwy$ for some $x \in \Sigma^*$ and $y \in \Sigma^+$.
	By assumption, $x = \varepsilon$ and $u = w$.
	However, as $u \notin L$ we have $u \ne w$.
	So we must have $(\pref{L} \setminus L)L \cap L \Sigma^+ = \emptyset$.
\end{proof}

\begin{corollary} \label{cor:solidcode-char2}
	Let $C \subseteq \Sigma^+$ be a solid code.
	Then $\pref{C}C \cap \Sigma^+ C \Sigma^+ = \emptyset$.
\end{corollary}

The condition $\pref{C}C \cap \Sigma^+ C \Sigma^+ = \emptyset$
clearly implies that~$C$ is comma-free,
but not every comma-free code fulfills it.
For example, for the comma-free code $C = \{aba,cd,bc\}$
we have $abcd \in (\pref{C}\setminus C)C \cap \Sigma bc \Sigma$.

The next examples show that $\pref{C}C \cap \Sigma^+ C \Sigma^+ = \emptyset$
does not characterize solid codes.

\begin{example}
	\begin{enumerate}
		\item For the comma-free code $C = \{abb,aba\}$ from \crefitem{constrsync:ex:binarysolidcode}{constrsync:ex:comma-free-not-solid}
		that is not a solid code we have $\pref{C}C \cap \Sigma^+ C \Sigma^+ = \emptyset$.
		
		\item Let $C = \{ab,bc\}$. Then $C$ is not a solid code,
		but we have $\pref{C}C \cap \Sigma^+ C \Sigma^+ = \emptyset$.
	\end{enumerate} 	
\end{example}

\begin{proposition} \label{constrsync:prop:prefLintersect}
	Let $L \subseteq \Sigma^+$.
	Then $\pref{L}L \cap \Sigma^+ L \Sigma^+ = \emptyset$
	if and only if for all $u \in \pref{L}$ and $v \in L$
	such that $uv = xwy$ for $x \in \Sigma^+$ and $y \in \Sigma^*$ and $w \in L$
	we have $y = \varepsilon$ and $w = v$.
\end{proposition}
\begin{proof}
	First, suppose $\pref{L}L \cap \Sigma^+ L \Sigma^+ = \emptyset$.
	In this case, we must have $y = \varepsilon$ since
	otherwise $uv \in \Sigma^+ L \Sigma^+$.
	Since~$L$ is a suffix code, this implies $w = v$.
	
	Conversely, the condition 
	implies $\pref{L}L \cap \Sigma^+ L \Sigma^+ = \emptyset$.	
\end{proof}

\begin{corollary} \label{cor:solidcode-char3}
	Let $L \subseteq \Sigma^+$.
	Then $L$ is a solid code
	if and only if for all $u \in \pref{L}$ and $v \in L$,
	if $uv = xwy$ for  $w \in L$ and $x,y \in \Sigma^*$,
	then $x = \varepsilon$ and $w = u$
	or $y = \varepsilon$ and $w = v$.
\end{corollary}
\begin{proof}
	\Cref{thm:solidcode-char2,cor:solidcode-char2,constrsync:prop:prefLintersect} show 
	that a solid code fulfills the condition.
	Conversely, it implies the condition
	from \cref{thm:solidcode-char2} if $|y| > 0$
	and thus that~$L$ is a solid code.
\end{proof}

Lastly, we give a characterization using the Levenshtein-Romanov mapping. It basically combines results
from \cref{constrsync:sec:codes:levensthein-romanov}.

\begin{proposition} \label{constrsync:prop:sigma-solid-code-char}
	Let $L \subseteq \Sigma^+$ be non-empty.
	Then $L$ is a solid code if and only if
	for all $u,v \in \Sigma^*$ the following is true:
	\begin{enumerate}
		\item \label{constrsync:prop:sigma-solid-code-char:1} 
		$\sigma_L(uv) = uv$ implies $u \notin L$ or $v = \varepsilon$,
		
		\item \label{constrsync:prop:sigma-solid-code-char:2} $\sigma_L(uv) \in L$ implies $v \notin \pref{L}\setminus L$ or $v = \varepsilon$,
		
		\item \label{constrsync:prop:sigma-solid-code-char:3} $\sigma_L(uv) \ne v$ implies $v \notin L$. 	
	\end{enumerate}
\end{proposition}
\begin{proof}
	First, suppose~$L$ is a solid code.	
	Assume $\sigma_L(uv) = uv$.
	Since~$L$ is a prefix code,
	by \cref{constrsync:prop:prefix-code:4} of \cref{constrsync:prop:prefix-code}
	we find that~$u \notin L$
	or $v = \varepsilon$.
	Next, let $\sigma_L(uv) \in L$ and $v \in \pref{L}\setminus L$.
	Then $\sigma_L(uv) = wv$ for some $w \in \Sigma^*$
	with $|w| > 0$ since $wv \in L$.
	By overlap-freeness, this yields $v = \varepsilon$.	
	\Cref{constrsync:prop:sigma-solid-code-char:3}
	follows from \cref{constrsync:prop:infix-code:6} of \cref{constrsync:prop:infix-code}
	since~$L$ is an infix code.
		
	Conversely, \cref{constrsync:prop:sigma-solid-code-char:1} 
	yields that~$L$ is a prefix code by \crefitem{constrsync:prop:prefix-code}{constrsync:prop:prefix-code:4}.
	Then, \cref{constrsync:prop:sigma-solid-code-char:3} 
	yields that~$L$ is an infix code by \cref{constrsync:prop:infix-code:6} 
	of \cref{constrsync:prop:infix-code}.
	Next, assume $uv, vw \in L$ for non-empty $u,v,w \in \Sigma^+$.
	We have $v \in \pref{L}\setminus L$ since~$L$ is a prefix code
	and $\sigma_L(uv) = uv$. Thus, \cref{constrsync:prop:sigma-solid-code-char:2} 
	does not hold. So, if all these conditions hold,
	we find that~$L$ is overlap-free.
\end{proof}

\begin{remark}
	By \cref{constrsync:prop:sigma-solid-code-char:3} 
	from \cref{constrsync:prop:sigma-solid-code-char},
	for a solid code~$C$ and $v \in C$,
	we have $\sigma_C(uv) = v$ for $u \in \Sigma^*$.
	By \cref{constrsync:prop:sigma-solid-code-char:2},
	we also have $\sigma_C(uw) \notin L$
	for every proper and non-empty prefix $w$
	of a word from~$C$.
	However, note that we do not have $\sigma_C(uw) = w$
	in this case. For example, for the solid code $C = \{abb, aac\}$
	we have $\sigma_C(uw) = uw$ for $u = a$
	and $w = a$.
\end{remark}